%% kapa: exam=1 round=32 year=2021 period=1 form=A subject=civil_law subject_ko=민법 seq=14 number=14
%% source: Q-Net 감정평가사 기출문제 (cst00309) artlSeq=5208397 file=extracted/1008/5208397/A형/2021년도 제32회 감정평가사 1차 1교시 A형.pdf
%% conversion: kordoc 4.24.51 -> tools/gen1.py (auto); review=proofread-sonnet; flags=none
\begin{question}{14}
의사표시의 효력발생에 관한 설명으로 옳지 않은 것은? (다툼이 있으면 판례에 따름)
\begin{choices}
  \item 도달주의의 원칙은 채권양도의 통지와 같은 준법률행위에도 유추적용될 수 있다.
  \item 의사표시의 부도달 또는 연착으로 인한 불이익은 특별한 사정이 없는 한 표의자가 이를 부담한다.
  \item 의사표시자가 그 통지를 발송한 후 제한능력자가 되었다면 특별한 사정이 없는 한 그 의사표시는 취소할 수 있다.
  \item 수령무능력자에게 의사표시를 한 경우, 특별한 사정이 없는 한 표의자는 그 의사표시로써 수령무능력자에게 대항할 수 없다.
  \item 상대방이 정당한 사유 없이 의사표시 통지의 수령을 거절한 경우, 상대방이 그 통지의 내용을 알 수 있는 객관적 상태에 놓여 있는 때에 의사표시의 효력이 생기는 것으로 보아야 한다.
\end{choices}
\end{question}
\answer{3}
